\chapter{Introducción}
\label{cap:introduccion}

\section{Motivación y antecedentes} 

\aer vio la luz por primera vez en 2014 de la mano de los profesores Marco Antonio y Pedro Pablo Gómez Martín de la Universidad Complutense de Madrid. Como podemos leer en \cite{Pedro_Pablo_Gomez_Martin_Marco_Antonio_Gomez_Martin_2017} este juez \textit{online} se creó como una herramienta educativa destinada a facilitar el aprendizaje de la programación, ofreciendo a los estudiantes un espacio en el que poner a prueba sus habilidades y conocimientos gracias a su amplio archivo de ejercicios en castellano. En los siguientes años siguió ganando popularidad, habiendo recibido ya más de un millón de envíos y contando con más de 5 mil usuarios\footnote{Puedes encontrar más hitos importantes de \aer aquí: \url{https://aceptaelreto.com/doc/history.php}}. Su utilidad como juez \textit{online} para ofrecer ejercicios interesantes en castellano ha provocado que su uso se haya extendido por centros educativos y haya servido de plataforma para concursos\footnote{\url{https://las12uvas.es/2025/\#/quees}}. Sin embargo, a pesar de su popularidad y su indiscutible valor pedagógico, \aer presenta limitaciones significativas en cuanto a funcionalidades. La plataforma carece de herramientas avanzadas para la gestión de clases o concursos, y su diseño minimalista no permite a los usuarios compartir o visualizar el código de otros participantes, lo que limita las posibilidades de colaboración y aprendizaje entre pares.

Otra de las carencias más notables es la ausencia de un sistema que permita a los usuarios evaluar su progreso de manera tangible, ya que las estadísticas que proporciona a los usuarios son limitadas. Aunque la plataforma ofrece una amplia gama de ejercicios organizados por categorías, facilitando la búsqueda de temas específicos según los intereses o necesidades de cada usuario \footnote{\url{https://aceptaelreto.com/problems/categories.php}} no incluye mecanismos para medir el avance individual, compararse con otros miembros de la comunidad o recibir reconocimientos por los logros alcanzados. 

Esta situación representa una oportunidad única para mejorar y enriquecer el servicio que ofrece \aer. Se puede observar que muchos de los jueces \textit{online} comerciales, especialmente aquellos enfocados en la mejora de habilidades técnicas (ejemplos: HackerRank, LeetCode), han logrado un éxito notable en la retención de usuarios gracias a la implementación de sistemas de puntos, recompensas y medallas. También podemos comprobar que existen páginas de terceros que pretenden ofrecer un servicio de estadísticas\footnote{\url{https://aer.lluiscab.net/}} para el juez, demostrando una demanda por parte de los usuarios, pero que al no ser una parte nativa de la aplicación no acceden directamente a la base de datos si no que realizan \textit{scraping}\footnote{Véase la cuarta definición para el término en el diccionario Cambridge: \url{https://dictionary.cambridge.org/dictionary/english/scraping}}. Esto genera el riesgo de que los datos se queden desactualizados al cambiar la UI de la aplicación. Además, solo actualizan los datos una vez cada cierta cantidad de tiempo\footnote{\url{https://aer.lluiscab.net/faq}} en lugar de actualizarse constantemente, por lo que a menudo están desactualizados y su utilidad para mostrar una imágen del progreso del usuario está limitada. \aer también cuenta con algunas estadísticas de usuario que se muestran en la página de perfil, pero éstas están limitadas, no hay estadísticas para ejercicios, ni se hace notar un sistema consistente de gamificación a lo largo de la aplicación.

En este sentido, \aer ya cuenta con una base sólida sobre la que integrar estas mejoras de manera nativa. La plataforma dispone de funcionalidades esenciales, como el seguimiento detallado de los envíos realizados por los usuarios, una categorización clara y accesible de los ejercicios, y una API robusta que permite el acceso a los datos de los usuarios de manera eficiente\footnote{\url{https://aceptaelreto.com/statistics/submissions.php}}. Estas características facilitan la implementación de un sistema de recompensas y visualización de progreso, que podría transformar la experiencia del usuario y elevar el valor educativo de la plataforma.

La motivación principal detrás de esta propuesta de mejora es, por tanto, añadir un valor significativo a la aplicación, reforzando su objetivo como espacio para la mejora de habilidades. La introducción de un sistema de medallas, junto con herramientas que permitan visualizar el progreso a lo largo del tiempo, no solo podría aumentar la motivación extrínseca de los usuarios, sino que también facilitaría la integración de la plataforma en entornos docentes (ejemplo: un profesor podría decidir otorgar un punto extra al alumno que obtenga más medallas durante un cierto periodo). 

\section{Objetivos}
Como objetivos tenemos diseñar un panel de estadísticas y un sistema de logros para el juez \textit{online} \aer, que permita mostrar información relevante de la plataforma, y que mejore la experiencia del usuario, fomentando su participación mediante técnicas de gamificación. A continuación se listan las metas específicas que queremos alcanzar con este proyecto:
\begin{enumerate}
    \item Identificar la información disponible en \aer que pueda utilizarse para la elaboración de estadísticas
    \item Establecer las métricas relevantes para representar la actividad de la plataforma.
    \item Diseñar un panel de estadísticas que permita visualizar de forma clara la información del juez.
    \item Mostrar estadísticas globales y de usuario que reflejen la actividad de la plataforma
    \item Determinar qué información y logros resultan relevantes para los usuarios del juez \textit{online}.
    \item Diseñar un sistema de logros e insignias asociado a la actividad de los usuarios.
    \item Mejorar la experiencia de usuario y la fidelización mediante elementos de gamificación.
\end{enumerate}


\section{Plan de trabajo}
Para alcanzar los objetivos descritos en el apartado anterior, el desarrollo del proyecto se dividirá en distintas fases. Esto nos permitirá abordar el proyecto de una forma estructurada. Seguiremos las metodologías de diseño de software gamificado \cite{Morschheuser_Hassan_Werder_Hamari_2018} que proponen un enfoque de diseño centrado en el usuario. En esta sección se describirá el proceso que seguiremos desde el análisis inicial del contexto hasta la implementación y evaluación de la solución final.

En la primera fase, se estudiará el funcionamiento del juez \textit{online} \aer, así como la información y los datos de los que dispone la plataforma. El objetivo de esta fase es identificar aquellos elementos del juez que puedan ser utilizados para la obtención de estadísticas relevantes. También, se analizarán plataformas similares y otros jueces \textit{online} para  identificar buenas prácticas en aspectos como la presentación de estadísticas o la implementación de técnicas de gamificación. Este análisis permitirá tener una visión global del sistema, identificando sus capacidades y restricciones de cara al proyecto.

En esta segunda fase, a partir de la información identificada en la fase anterior, se establecerán las métricas que permiten reflejar la actividad del juez \textit{online}, tanto a nivel global como a nivel de usuario. Estas métricas se seleccionarán teniendo en cuenta su relevancia para los usuarios y su utilidad para evaluar la actividad y evolución de la plataforma. 

Durante la tercera fase, se diseñará el panel de estadísticas y los distintos \textit{dashboards} que lo componen que permitirán mostrar tanto las estadísticas globales de la plataforma como las estadísticas individuales de los usuarios. Para ello, se estudiarán diferentes tipos de paneles y técnicas de visualización de la información, con el objetivo de seleccionar aquellas que permitan representar las métricas de forma clara y comprensible. 

En la cuarta fase, se analizarán las necesidades, motivaciones y perfiles de los usuarios del juez \textit{online}, apoyándose en un cuestionario, y de esa forma, diseñar un sistema de gamificación adaptado a esas necesidades. A partir de este análisis, y siguiendo los distintos principios de diseño de software gamificado, se definirán los diferentes elementos de gamificación ya sean puntos, niveles o insignias \cite{Zichermann_Cunningham_2011}. Además, se diseñará el sistema de logros asociado a la actividad del usuario, para fomente su participación.

La quinta fase consistirá en implementar el panel de estadísticas y el sistema de gamificación diseñados en las fases anteriores. Tal y como se indica en libro mencionado anteriormente, la implementación se realizará siguiendo un enfoque iterativo, lo que permitirá realizar ajustes progresivos en función de los resultados obtenidos y de las pruebas realizadas. 

Finalmente, en la sexta fase, se realizarán pruebas para verificar el correcto funcionamiento del sistema desarrollado y comprobar que se cumplen los objetivos planteados. 